\documentclass[11pt,a4paper]{article}

\usepackage[utf8]{inputenc}

\usepackage[T1]{fontenc}

\usepackage[polish]{babel}

\usepackage{lmodern}

\usepackage{geometry}

\geometry{a4paper,margin=27mm}

\usepackage{microtype}

\usepackage[hidelinks]{hyperref}

% Portfolio: MSC-AR-P04; source: prospective ARC-XR research scope.

\title{Grounding przestrzenny obiektów w instrukcjach AR na podstawie VLM}

\author{Mateusz Pomianek}

\date{Wersja 0.4.0, 26 listopada 2026}

\begin{document}

\maketitle

\begin{abstract}

Przedmiotem opracowania jest Grounding przestrzenny obiektów w instrukcjach AR na podstawie VLM. Czy dołączenie kalibrowanej geometrii 3D do wejścia VLM poprawia wskazywanie obiektu względem modelu 2D? Ten sam model VLM zostanie oceniony na jednakowych scenach w wariancie obrazowym 2D i w wariancie z niezależnie skalibrowaną geometrią 3D. Głównym punktem końcowym jest odsetek prawidłowo wskazanych obiektów na niezależnym zbiorze scen. Opis ma charakter wstępny i nie przedstawia niezweryfikowanych wyników.

\end{abstract}

\section{Wprowadzenie}

Rozwój adaptacyjnych systemów XR stwarza możliwość badania interakcji człowieka z informacją przestrzenną, ale samo zastosowanie AR, VR lub modeli sztucznej inteligencji nie przesądza o przewadze poznawczej czy ergonomicznej. W rozpatrywanym zagadnieniu, Grounding przestrzenny obiektów w instrukcjach AR na podstawie VLM, kluczowe jest rozdzielenie właściwości systemu od zmienności zachowania użytkownika. Porównanie należy zaprojektować tak, aby mierzyć jednoznaczne zjawisko, a nie tylko ogólną ocenę technologii.

\section{Cel pracy i pytanie badawcze}

Czy dołączenie kalibrowanej geometrii 3D do wejścia VLM poprawia wskazywanie obiektu względem modelu 2D? Wariant minimalny nie wymaga bezpośredniej akwizycji głębi z HMD: niezależna referencja przestrzenna może pochodzić z kalibrowanej sceny.

\section{Zakres i odrębny wkład naukowy}
Przestrzenne wskazanie poprawnego obiektu referencyjnego przez ten sam VLM w wersji obrazowej 2D i z dostarczoną, niezależnie kalibrowaną geometrią 3D. Test przyrostu poprawności kotwic w określonej scenie przy identycznym VLM, nie deklaracja wynalezienia grounding 3D.

\section{Hipotezy badawcze}
\paragraph{H1.} Dodanie niezależnie kalibrowanych danych 3D zwiększa odsetek poprawnych wskazań względem tego samego VLM wyłącznie obrazowego.
\paragraph{H2.} Zysk jest większy dla scen z dwoma podobnymi obiektami i relacjami przestrzennymi niż dla jednoznacznych scen pojedynczego celu.
\paragraph{H3.} Korzyść utrzymuje się na niewidzianych układach scen bez wyboru najlepszego promptu po analizie zbioru testowego.

\section{Projekt badania i plan eksperymentów}
\subsection{A. Operacjonalizacja i kwalifikacja danych}
Przestrzenne wskazanie poprawnego obiektu referencyjnego przez ten sam VLM w wersji obrazowej 2D i z dostarczoną, niezależnie kalibrowaną geometrią 3D. Jednostka wnioskowania: scena z jednoznaczną referencją obiektu; powtórzenia promptu zagnieżdżone w scenie. Punkt końcowy i sposób oceny będą zamrożone przed próbą potwierdzającą. Należy udokumentować pochodzenie obserwacji, czas akwizycji, identyfikatory sesji, harmonogram oraz legalność użycia materiałów.

\subsection{B. Eksperyment porównawczy}
Zamrozić VLM, prompty, listę kandydatów, obrazy i reguły dopasowania; zmieniać jedynie obecność geometrii 3D. Estymand pierwszorzędowy: różnica odsetka poprawnych referencji obiektowych między VLM+3D a tym samym VLM 2D. Przypisanie warunków, mianowniki i kryteria wyłączeń nie mogą być dobierane po przejrzeniu wyników.

\subsection{C. Kontrole i analiza odporności}
Niezależnie pomierzyć kotwiczenie i współrzędne referencyjne, nie traktować samego syntetycznego pokrycia 2D jako ground truth. Okluzje, niejednoznaczne nazwy i podobne obiekty; raportować odmowy wskazania i błędne identyfikacje. Każdą odmowę systemu, niekompletną próbę i brak sygnału należy raportować osobno, a nie po cichu usuwać z próby.

\subsection{D. Walidacja na niezależnym materiale}
Ocenę należy wykonać na uprzednio odłożonych scenach, sesjach lub uczestnikach. Parowany model Bernoulliego z efektem sceny, podziałem po geometrii scen i przedziałem różnicy trafności. Model i próg zatrzymać przed analizą końcową; rejestrować również brak przewagi i działania niepożądane.

\section{Estymand, analiza statystyczna i decyzje}
Główny estymand: różnica odsetka poprawnych referencji obiektowych między VLM+3D a tym samym VLM 2D. Jednostką powtarzalności jest scena z jednoznaczną referencją obiektu; powtórzenia promptu zagnieżdżone w scenie. Parowany model Bernoulliego z efektem sceny, podziałem po geometrii scen i przedziałem różnicy trafności. H1 jest kandydatem na hipotezę potwierdzającą. H2 i H3 pozostają eksploracyjne, dopóki protokół nie określi ich hierarchii i kontroli wielokrotnych porównań. Minimalny efekt istotny praktycznie, liczebność prób, margines nie gorszości i progi jakości wyznaczyć na podstawie pilotażu i analizy precyzji, nie zgadywać wartości.

\section{Stan badań i granica nowości}
\begin{itemize}
\item 3D grounding with VLM is prior art; the new study must demonstrate calibrated AR anchor gains, not merely combining 2D and 3D. \cite{xrprior1}.
\item Task-adaptive spatial language retrieval and meter-accurate anchors are already proposed, narrowing independent novelty. \cite{xrprior2}.
\end{itemize}
Test przyrostu poprawności kotwic w określonej scenie przy identycznym VLM, nie deklaracja wynalezienia grounding 3D. Samo zastosowanie VLM, EEG, okulometrii lub adaptacyjnego HUD nie stanowi wystarczającej nowości międzynarodowej.

\section{Warunki bezpieczeństwa, sprzętu i zgodności}
Opis stanowi prospektywny plan, a nie prerejestrację ani zatwierdzoną procedurę. Przed zbieraniem danych od ludzi niezbędna jest właściwa weryfikacja etyczna i zgoda. Standardowy Quest 3 nie ma wbudowanego eye trackingu, a Tobii Eye Tracker 5 działa na stanowisku ekranowym. Odstępstwa, obciążenia i ograniczenia czujników muszą być ujawnione. Nie publikować surowych danych biometrycznych ani nagrań uczestników.

\section{Materiał i metody}

Ten sam model VLM zostanie oceniony na jednakowych scenach w wariancie obrazowym 2D i w wariancie z niezależnie skalibrowaną geometrią 3D. Wymagane jest zdefiniowanie jednostki analizy, zasad doboru przypadków, wersji scenariusza i źródeł pomiaru. Rozkład warunków oraz reguły wykluczenia prób należy określić przed pomiarem. W badaniach z udziałem ludzi niezbędna jest uprzednia właściwa ocena etyczna i świadoma zgoda. W testach wyłącznie symulacyjnych należy oddzielić prawdę symulacyjną od obserwacji empirycznej.

\subsection{Wynik główny i kryteria jakości}

Odsetek prawidłowo wskazanych obiektów na niezależnym zbiorze scen. Pomiar powinien zachowywać jawne jednostki, znaczniki czasu i reguły oceny poprawności. Przy modelach predykcyjnych stosować rozłączne partycje oceny; przy eksperymentach interwencyjnych oddzielić skuteczność predykcji od przyczynowego wpływu interfejsu.


\subsection{Benchmark syntetyczny i metody referencyjne}
W prototypie technicznym zastosowano kamerę otworkową o rozdzielczości
$640\times480$ pikseli i ogniskowej $f_x=f_y=440$ pikseli.
Dla punktu $(X,Y,Z)$, gdzie $Z>0$, współrzędne obrazu wynoszą
\[
 u=c_x+f_x X/Z,\qquad v=c_y-f_y Y/Z.
\]
Sceny zawierają trzy oznaczone prostopadłościany o zróżnicowanej wysokości
i głębokości. Generator zapisuje obraz RGB, geometrię oraz oddzielne
referencje odpowiedzi. W wariancie podstawowym przewidziano 30 scen,
z których 10 służy rozwojowi, a 20 pozostaje odłożonych do późniejszej oceny.
Porównanie główne obejmuje ten sam VLM z obrazem RGB oraz z obrazem
i jawną geometrią. Dodatkowy algorytm deterministyczny wybiera obiekt
o najmniejszej głębokości lub największej wysokości, korzystając z identyfikatora
szablonu instrukcji. Nie stanowi zatem równoważnego testu rozumienia języka.
Jego czas obliczeń nie jest bezpośrednio porównywalny z pełnym czasem
żądania do serwera VLM.

\subsection{Eksploracyjna analiza odporności}
Proponowana macierz obejmuje odchylenia standardowe szumu współrzędnych
$0$, $0{,}02$ i $0{,}05$ m oraz zasłonięcie $0$, $10$ i $25$ procent
szerokości obrazu pionowym pasem. Zakłócenia nie zmieniają referencji
odpowiedzi. Jest to uproszczona maska obrazu, nie model fizycznej okluzji.
Dla sceny $i$ z $m_i$ instrukcjami definiuje się
\[
 d_i(g,o)=\frac{1}{m_i}\sum_{j=1}^{m_i}
 (C_{ij}^{3D}(g,o)-C_{ij}^{RGB}(g,o)),
\]
gdzie $C$ jest wskaźnikiem poprawności. Estymator przewagi geometrii
$\widehat\Delta_{g,o}=n^{-1}\sum_{i=1}^{n}d_i(g,o)$ nadaje scenom jednakowe wagi.
Eksploracyjny kontrast interakcji wynosi
\[
 \widehat I_{g,o}=\widehat\Delta_{g,o}-\widehat\Delta_{g,0}
 -\widehat\Delta_{0,o}+\widehat\Delta_{0,0}.
\]
Bootstrap percentylowy wykorzystuje 2000 replikacji całych scen, zachowując
parowanie pomiędzy konfiguracjami. Przedziały są punktowe, bez korekty
wielokrotnych porównań. Przy dwóch scenach pilotażowych służą wyłącznie
sprawdzeniu implementacji. Odmowa lub niepoprawny format odpowiedzi są
liczone jako niepoprawne wskazanie; błąd wykonania blokuje kompletną
analizę macierzy. Pilotaż rozwojowy nie rozstrzyga H1--H3.

\section{Plan analizy}

Porównać dokładność na poziomie sceny i obiektu, analizować klasy pomyłek, abstencję i wpływ błędu kalibracji. W wynikach należy raportować wielkość efektu lub błąd estymacji z przedziałem niepewności, liczbą obserwacji i odsetkiem odrzuconych prób. Analizy dodatkowe wymagają jasnego oznaczenia jako eksploracyjne. Nie wybierać progów decyzyjnych ani modelu na podstawie obserwacji ze zbioru testowego.

\section{Wyniki}

W obecnej wersji nie raportuje się wyników pomiarów, tabel statystycznych ani parametrów wyznaczonych z nieudokumentowanych danych. Nie stwierdza się w tym opracowaniu, że eksperyment został wykonany lub zgłoszony do czasopisma.

\section{Dyskusja}

Interpretacja przyszłych wyników musi uwzględnić alternatywne mechanizmy i zależność wniosków od jakości danych, warunków zadania oraz ograniczeń sprzętu. Wartość potencjalnego rozwiązania zależy nie tylko od trafności modelu, lecz także od jego opóźnienia, stabilności i wpływu na zachowanie człowieka. Wyniki testów syntetycznych nie mogą być automatycznie uogólniane na populację uczestników.

\section{Ograniczenia badania}

Nie zakładać dostępu do map głębi Quest 3; geometria musi pochodzić z udokumentowanego, kalibrowanego źródła. Prawidłowe przypisanie przestrzenne, odrębne od weryfikacji kolejności kroków P05. Żadna część szkicu nie dowodzi przeprowadzenia zbierania danych, uzyskania zgód etycznych ani gotowości manuskryptu do zgłoszenia.

\section{Wnioski}

Przedstawiono problem, kontrast i sposób jego planowanej oceny. Rozstrzygnięcie postawionego pytania wymaga przygotowania protokołu, zweryfikowanych źródeł oraz poprawnie udokumentowanych wyników. Bez tych materiałów nie można sformułować konkluzji empirycznych.

\begin{thebibliography}{9}
\bibitem{xrprior1} SeeGround: See and Ground for Zero-Shot Open-Vocabulary 3D Visual Grounding (2025), CVPR 2025. \url{https://ieeexplore.ieee.org/abstract/document/11094047}.
\bibitem{xrprior2} Words into World: A Task-Adaptive Agent for Language-Guided Spatial Retrieval in AR (2025), arXiv:2512.00294. \url{https://arxiv.org/abs/2512.00294}.
\end{thebibliography}
\end{document}

